% from ega3-1-fr.tex
\oldpage[III]{146}

\noindent\textit{II)} Revenons aux notations du début de la démonstration,
et remarquons qu'on peut toujours supposer $Y=\operatorname{Spec}(A)$
affine~; comme $f'$ est de type fini et $\mathcal L'$ ample pour $f'$, il
existe un entier $m>0$ tel que $\mathcal L'^{\otimes m}$ soit très ample pour
$f'$ (II, 4.6.11)~; remplaçant au besoin $\mathcal L$ par
$\mathcal L^{\otimes m}$, on peut se borner à considérer le cas où
$\mathcal L'$ est très ample pour $f'$, et à prouver que
$\mathcal L|f^{-1}(U)$ est alors très ample pour $f$. Comme $f'$ est propre,
il existe alors une $Y'$-immersion fermée
$j:X'\to P=\mathbf P_{Y'}^r$ pour un entier $r>0$ convenable, telle que
$\mathcal L'$ soit isomorphe à $j^*(\mathcal O_P(1))$ (II, 5.5.4, (ii))~;
cette immersion correspond canoniquement à un $\mathcal O_{X'}$-homomorphisme
surjectif
\[
u:\mathcal O_{X'}^{r+1}\longrightarrow\mathcal L'
\]
(II, 4.2.3). Ce dernier correspond (0\textsubscript{I}, 5.1.1) à la donnée
de $r+1$ sections $s'_i$ ($0\leq i\leq r$) de $\mathcal L'$ au-dessus de
$X'$ qui engendrent ce $\mathcal O_{X'}$-Module. Ces sections sont aussi par
définition des sections de $f'_*(\mathcal L')$ au-dessus de $Y'$~; on a
\[
f'_*(\mathcal L')=f_*(\mathcal L)\otimes_{\mathcal O_Y}\mathcal O_{Y'}
\]
(0\textsubscript{I}, 5.4.10), $Y$ est affine et $\mathcal O_y$ est l'anneau
local en l'idéal premier $\mathfrak j_y$ de $A$, donc on a
$s'_i=s''_i/t_i$, où les $s''_i$ sont des sections de $f_*(\mathcal L)$
au-dessus de $Y$ et les $t_i$ des éléments de $A$ n'appartenant pas à
$\mathfrak j_y$~; on en conclut qu'il existe un voisinage ouvert affine $V$
de $y$ dans $Y$ et des sections $s_i$ de $f_*(\mathcal L)|V$ telles que
$s'_i=s_i/1$ (on rappelle que l'espace $Y'$ est contenu dans $V$, cf. I,
2.4.2). Les $s_i$ sont alors des sections de $\mathcal L$ au-dessus de
$f^{-1}(V)$, définissant donc un homomorphisme
\[
v:\bigl(\mathcal O_X|f^{-1}(V)\bigr)^{r+1}
\longrightarrow\mathcal L|f^{-1}(V)
\]
qui, par hypothèse, est surjectif en tous les points de $f^{-1}(y)$~; comme
$\operatorname{Coker}(v)$ est cohérent (0\textsubscript{I}, 5.3.4), son
support est fermé (0\textsubscript{I}, 5.2.2) et par suite il existe un
voisinage ouvert $W\subset f^{-1}(V)$ de $f^{-1}(y)$ tel que la restriction
de $v$ à $W$ soit un homomorphisme surjectif. Puisque le morphisme $f$ est
fermé, on peut supposer que $W$ est de la forme $f^{-1}(U)$, où $U$ est un
voisinage ouvert de $y$, et la conclusion résulte alors de (II, 4.2.3).
\end{proof}

\subsection{Morphismes finis de préschémas formels}

\begin{proposition}[4.8.1]
\label{III.4.8.1-fr}
Soient $\mathfrak Y$ un préschéma formel localement noethérien, $\mathcal K$ un
Idéal de définition de $\mathfrak Y$, $f:\mathfrak X\to\mathfrak Y$ un
morphisme de préschémas formels. Les conditions suivantes sont équivalentes~:
\begin{enumerate}
\item[a)] $\mathfrak X$ est localement noethérien, $f$ est un morphisme adique
  (I, 10.12.1) et si l'on pose
  $\mathcal J=f^*(\mathcal K)\mathcal O_{\mathfrak X}$, le morphisme
  \[
  f_0:(\mathfrak X,\mathcal O_{\mathfrak X}/\mathcal J)
  \longrightarrow
  (\mathfrak Y,\mathcal O_{\mathfrak Y}/\mathcal K)
  \]
  déduit de $f$ est fini.
\item[b)] $\mathfrak X$ est localement noethérien et est limite inductive d'un
  $(Y_n)$-système inductif adique $(X_n)$ tel que le morphisme
  $X_0\to Y_0$ soit fini.
\item[c)] Tout point de $\mathfrak Y$ possède un voisinage ouvert formel affine
  noethérien $V$ tel que $f^{-1}(V)$ soit un ouvert formel affine et que
  $\Gamma(f^{-1}(V),\mathcal O_{\mathfrak X})$ soit un
  $\Gamma(V,\mathcal O_{\mathfrak Y})$-module de type fini.
\end{enumerate}
\end{proposition}

\begin{proof}
Il est immédiat que a) entraîne b) en vertu de (I, 10.12.3). Pour voir que b)
entraîne c), on peut supposer que
$\mathfrak Y=\operatorname{Spf}(B)$, où $B$ est adique noethérien et
$\mathcal K=\mathfrak R^\Delta$, où $\mathfrak R$ est un idéal de définition
de $B$. Par hypothèse, $X_0$ est un schéma affine dont l'anneau $A_0$ est un
$B/\mathfrak R$-module de type fini (II, 6.1.3). En vertu de (I, 5.1.9), chacun
des $X_n$ est un schéma affine, et si $A_n$ est son anneau, l'hypothèse b)
entraîne que pour $m\leq n$, $A_m$ est isomorphe à
$A_n/\mathfrak R^{m+1}A_n$. On en déduit que $\mathfrak X$ est isomorphe à
$\operatorname{Spf}(A)$, où $A=\varprojlim A_n$~; on conclut en vertu de
(0\textsubscript{I}, 7.2.9). Enfin, pour prouver que c)

